%!TEX program = lualatex
\documentclass[a4paper,ngerman,headheight=80pt,12pt,DIV=calc]{scrartcl}
\usepackage[ngerman]{babel}
\usepackage{blindtext}
\usepackage[utf8]{luainputenc}
\usepackage{scrlayer-scrpage}
\usepackage[T1]{fontenc}
\usepackage{fontspec}
\usepackage{multicol}
\usepackage{enumitem}
\usepackage{epigraph}
\usepackage[german=quotes]{csquotes}
\usepackage[footskip=20pt,bottom=50pt]{geometry}
\setmainfont{Source Sans Pro}
\MakeOuterQuote{"}

\setlength{\parindent}{0pt}
\setlength{\parskip}{6pt}

\pagestyle{plain.scrheadings}
\thispagestyle{scrheadings}
\lohead{\vspace*{-1.1cm}\includegraphics[scale=0.3]{../images/stupa_logo_rot.pdf}}
%\rohead{\\Präsidium des\\Studierendenparlaments}

\lofoot[]{\small V.i.S.d.P. Gunhild Berdal, Jim Martens, Ramon Weilinger c/o Von-Melle-Park 5, 20146 Hamburg}
\cofoot[]{}
\cohead[\thepage]{}

\renewcommand{\thesubsection}{TOP \arabic{subsection}}
\renewcommand{\thesubsubsection}{\arabic{subsubsection}.}
\setcounter{subsection}{-1}

\begin{document}
    UHH · StuPa-Präsidium · Von-Melle-Park 5 · D-20146 Hamburg\\
    Tel.: 040/450 204 39 · stupa@uni-hamburg.de

    \subsection*{Die Briefwahl zum Studierendenparlament geht los}
    \subsubsection*{"Fremdbestimmt oder selbstgestaltend? Universität gemeinsam gestalten!"}

    So lautet das Motto der StuPa-Wahl für die Legislatur 2017/18. Lassen wir
    uns von fremden Interessen, weltfremden politischen Vorgaben und restriktiven
    Anforderungen "treiben", oder nehmen wir die Dinge in die Hand und gestalten
    selbst die Inhalte und Strukturen des wissenschaftlichen Lebens?
    Die weitere Entwicklung der Universitäten ist umstritten: Soll die Studienreform
    die Bildung mündiger Persönlichkeiten und gesellschaftlich verantwortlicher
    Wissenschaft vorantreiben? Wird das Studium durch eine Rahmenprüfungsordnung
    entschleunigt und eine gemeinsame Lernkultur ermöglicht? Gelingt es, eine
    bedarfsgerechte Finanzierung der Hochschulen durchzusetzen?
    Jede/r Einzelne ist gefragt, sich einzubringen und Einfluss zu nehmen!

    Diese Tage bekommst Du per Post die Unterlagen für die \textbf{Briefwahl} zum
    Studierendenparlament 2017/18 zugeschickt. Darin findest Du eine Wahlbroschüre
    mit Informationen zur Wahl und den Darstellungen der kandidierenden Listen.
    Wenn Du per Brief wählst, beachte bitte die Hinweise in der Wahlbroschüre.
    Spätestens am \textbf{2. Januar 2017} muss der Rückumschlag bei uns angekommen sein.
    Den Brief kannst Du auch in die Urne im AStA Info-Café (Von-Melle-Park 5)
    einwerfen.

    Vom \textbf{16. bis 20. Januar 2017} findet die \textbf{Urnenwahl} statt,
    hier kannst Du auch persönlich an einer Urne die Stimme für das
    Studierendenparlament abgeben. Es kann insgesamt nur EINE Stimme abgegeben
    werden – entweder per Brief oder per Urne.

    (Im Verlaufe des Dezembers erhältst Du auch die Wahlunterlagen zum
    Akademischen Senat. Diese Wahl findet als reine Briefwahl statt. Bitte achte
    auf die Inhalte und wirf nicht versehentlich Wahlunterlagen – und Deine
    Stimme – weg.)

    Damit Du eine qualifizierte Entscheidung treffen kannst, organisieren wir
    am \textbf{Donnerstag, den 15. Dezember} eine öffentliche Listenvorstellung,
    wo alle kandidierenden Listen Rede und Antwort stehen.

    Zuerst gibt es einen politischen \textbf{"Markt der Möglichkeiten"}, wo die Aktiven
    mit Info-Tischen präsent und ansprechbar sind, und zwar von \textbf{16-18 Uhr
    in der WiWi-Mensa}.

    Anschließend gibt es um \textbf{18 Uhr im Raum 0079 des WiWi-Bunkers}
    (Von-Melle-Park 5) eine \textbf{Debatte}, in der sich die Listen grundsätzlichen
    weltanschaulichen Fragen wie den Folgenden stellen:

    \begin{itemize}
        \item Wie soll Hamburg in zehn Jahren aussehen?
        \item Wie soll die Gesellschaft (dann) aussehen?
        \item Wie soll sich die Universität bis dahin entwickelt haben?
    \end{itemize}

    Alle sind eingeladen, Fragen zu stellen und mitzudiskutieren!

\end{document}
